\clearpage
\phantomsection

\addcontentsline{toc}{chapter}{{Mở đầu}}
\chapter*{Mở đầu}

\noindent{\Large \textbf{Bài toán}}

Game hành động Roguelike đặt ra ba thách thức đồng thời cho tác nhân học tăng cường. Thứ nhất, \textit{phân phối bản đồ thay đổi theo từng lượt chơi (episode)}: tác nhân không thể ghi nhớ bố cục cố định mà phải khái quát hóa chính sách qua nhiều cấu hình không gian khác nhau. Thứ hai, \textit{không gian hành động đa nhánh $(3,3,3,9)$}: tác nhân phải phối hợp đồng thời hướng di chuyển, quyết định chiến đấu và góc ngắm trong thời gian thực. Thứ ba, \textit{phần thưởng thưa và bị trì hoãn}: các tín hiệu có giá trị cao như tiêu diệt kẻ địch hoặc dọn sạch phòng chỉ xuất hiện sau khi tác nhân hoàn thành chuỗi hành vi liên tiếp gồm tiếp cận, né tránh, căn hướng và tấn công.

Sự kết hợp của ba thách thức này khiến các kỹ thuật đơn lẻ như phần thưởng nội tại (ICM, RND) hay bộ nhớ tuần tự (LSTM) chưa chắc tạo ra cải thiện rõ rệt, từ đó đặt ra câu hỏi: cần kết hợp cơ chế nào để tác nhân học được hành vi chiến đấu hiệu quả trên môi trường sinh theo thủ tục?

\vspace{0.3cm}
\noindent\textbf{Phát biểu bài toán.} Cho môi trường Unity sinh phòng chiến đấu Roguelike theo thủ tục ở mỗi lượt chơi, đề tài huấn luyện một chính sách $\pi_\theta(a|s)$ thỏa mãn:

\begin{itemize}
    \item \textbf{Đầu vào (quan sát $s_t$)}: vector cấu trúc mô tả trạng thái người chơi, kẻ địch trong tầm nhìn, đạn đang bay và tường xung quanh (cấu hình cơ sở: 207 chiều).
    \item \textbf{Đầu ra (hành động $a_t$)}: bốn nhánh rời rạc $(3,3,3,9)$ gồm di chuyển trục X, di chuyển trục Y, ý định chiến đấu (none/attack/dash) và hướng tấn công.
    \item \textbf{Hàm thưởng $r_t$}: kết hợp tín hiệu kết quả thưa (hạ kẻ địch, dọn phòng) và tín hiệu định hình dày (tiếp cận, căn hướng tấn công).
    \item \textbf{Mục tiêu tối ưu}: tìm chính sách $\pi_\theta$ cực đại hóa kỳ vọng tổng phần thưởng.
    \item \textbf{Tiêu chí thành công}: tỷ lệ dọn phòng tail-50 $\geq 0.5$ trên phân phối bản đồ Hard do PCGNN sinh.
\end{itemize}

\vspace{0.5cm}
\noindent{\Large \textbf{Mục tiêu đề tài}}

Đề tài xác định bốn mục tiêu cụ thể sau, là cơ sở để thiết kế thực nghiệm và đánh giá kết quả xuyên suốt các chương:

\renewcommand{\labelitemi}{$-$}
\begin{itemize}
    \item \textbf{MT1}: Xây dựng tác nhân RL có thể chiến đấu hiệu quả trong môi trường Roguelike sinh theo thủ tục, đánh giá qua \textit{clear\_rate} và \textit{kills/episode}.
    \item \textbf{MT2}: Giải quyết bài toán phần thưởng thưa và phân phối bản đồ thay đổi bằng cách kết hợp học theo chương trình (curriculum learning) với bản đồ do PCGNN sinh.
    \item \textbf{MT3}: Xác định kỹ thuật hỗ trợ nào (ICM, RND, LSTM) thực sự tác động đến hiệu năng thông qua ma trận thí nghiệm loại trừ.
    \item \textbf{MT4}: Kiểm tra giả thuyết rằng biểu diễn quan sát là nút thắt của bài toán, đồng thời đánh giá khả năng cải thiện điều hướng mà không cung cấp trực tiếp hướng đi tối ưu cho tác nhân.
\end{itemize}

\vspace{0.5cm}
\noindent{\Large \textbf{Giải pháp đề xuất}}

Để đạt được các mục tiêu này, đề tài xây dựng một quy trình huấn luyện hoàn chỉnh gồm bốn thành phần chính: (1) PCGNN sinh bản đồ và BFS kiểm tra tính hợp lệ; (2) bộ phân loại độ khó chia bản đồ thành Easy/Medium/Hard và TilemapGenerator nạp vào Unity theo curriculum learning; (3) CombatAgent với vector quan sát cấu trúc và không gian hành động đa nhánh huấn luyện bằng PPO; (4) mở rộng quan sát lên 432 chiều bằng lưới cục bộ $15\times15$ (local occupancy grid) trong run62. Ma trận thí nghiệm loại trừ gồm 9 lần chạy chính cho phép cô lập tác động của từng kỹ thuật; run59 và run62 kiểm tra giả thuyết nút thắt biểu diễn quan sát.

\textbf{Cải tiến từ bài báo PCGNN \cite{Beukman2022}.} Quy trình sinh bản đồ kế thừa NEAT + novelty search của Beukman và cộng sự, nhưng được điều chỉnh ba điểm để phục vụ bài toán chiến đấu Roguelike thay vì sinh mê cung: (i) bộ hậu xử lý đặt vị trí người chơi/kẻ địch (P/E) trên cùng thành phần liên thông và kiểm tra BFS để bảo đảm tương tác chiến đấu hợp lệ; (ii) thay nhóm chỉ số gốc (leniency, A* difficulty, compression distance) bằng \emph{reachable\_ratio}, \emph{dead\_end\_ratio} và \emph{difficulty\_score} -- ba chỉ số phản ánh trực tiếp không gian di chuyển và né đạn; (iii) tách generator chạy ngoại tuyến để tránh phụ thuộc Python--Unity lúc huấn luyện, đồng thời cho phép kiểm tra bản đồ trước khi đưa vào tập huấn luyện. Chi tiết các cải tiến này được trình bày ở Chương~3.

\vspace{0.5cm}
\noindent{\Large \textbf{Câu hỏi nghiên cứu}}

Từ các mục tiêu trên, đề tài xác định hai câu hỏi nghiên cứu chính:

\renewcommand{\labelitemi}{$-$}
\begin{itemize}
    \item \textbf{Q1}: Trong số ICM, RND, LSTM và curriculum learning, kỹ thuật nào tạo ra cải thiện lớn nhất khi huấn luyện tác nhân chiến đấu trong môi trường Roguelike có bản đồ sinh theo thủ tục?
    \item \textbf{Q2}: Khi kết hợp curriculum learning với các kỹ thuật trên, có xuất hiện hiệu ứng bổ trợ rõ rệt hay không?
\end{itemize}

Để trả lời hai câu hỏi này, đề tài triển khai ma trận thí nghiệm loại trừ gồm chín lần chạy chính, trong đó mỗi lần chạy thay đổi đúng một biến so với cấu hình tham chiếu. Mỗi cấu hình chỉ được chạy một seed; vì vậy các so sánh được hiểu là xu hướng thực nghiệm, chưa phải kết luận thống kê có khoảng tin cậy.

\vspace{0.3cm}

\noindent{\Large \textbf{Đóng góp của đề tài}}

Thông qua quá trình nghiên cứu, triển khai và thực nghiệm, đề tài đạt được các đóng góp chính sau:

\renewcommand{\labelitemi}{$-$}
\begin{itemize}
    \item Xây dựng quy trình huấn luyện DRL cho game Roguelike trong Unity kết hợp PCGNN, curriculum learning và PPO.
    \item \textbf{Cải tiến PCGNN cho bài toán chiến đấu} so với thiết lập gốc của Beukman \cite{Beukman2022}: bổ sung bộ hậu xử lý đặt P/E + kiểm tra BFS, thay nhóm chỉ số gốc bằng \emph{reachable\_ratio}/\emph{dead\_end\_ratio}/\emph{difficulty\_score} phù hợp với không gian chiến đấu, và tách generator ngoại tuyến để tăng độ ổn định huấn luyện.
    \item Thiết kế CombatAgent với quan sát 207 chiều và không gian hành động đa nhánh $(3,3,3,9)$, phù hợp cho bài toán di chuyển--ngắm--chiến đấu đồng thời.
    \item Bằng chứng thực nghiệm cho thấy curriculum learning cải thiện phần thưởng trung bình \textbf{+85\%} (2.073 → 3.843) so với cấu hình PPO cơ sở, trong khi ICM, RND và LSTM chỉ đạt +16\%--20\%.
    \item Phát hiện vùng bão hòa entropy 4.40 nat nhất quán trên chín cấu hình thí nghiệm loại trừ, cho thấy đây là hệ quả của quan sát thiếu đặc trưng không gian chứ không phải giới hạn của PPO hay không gian hành động.
    \item Run59 (\emph{path-informed}, bổ sung đặc trưng BFS) cho thấy tìm đường là một nút thắt quan trọng; run62 (432 chiều, 40.82M bước) đưa entropy xuống 3.048 nat và đạt clear\_rate 0.666 dù không được cung cấp trực tiếp hướng đi tối ưu.
\end{itemize}

\vspace{0.3cm}

\noindent{\Large \textbf{Bố cục của dự án}}

\renewcommand{\labelitemi}{$-$}
\begin{itemize}
	\item \textbf{Mở đầu}: Bài toán, bốn mục tiêu của đề tài, giải pháp tổng quan và câu hỏi nghiên cứu.
	\item \textbf{Chương 1}: Cơ sở lý thuyết -- RL, PPO, GAE và các kỹ thuật hỗ trợ (ICM, RND, LSTM, PCGNN, curriculum), phân tích hạn chế và cách đề tài xử lý. Giới thiệu đặc trưng của Roguelike làm nền tảng cho thiết kế.
	\item \textbf{Chương 2}: Thiết kế hệ thống -- Tổng quan quy trình đầu cuối, đặc tả POMDP, quan sát, không gian hành động, hàm phần thưởng, quy trình PCGNN và khung curriculum learning.
	\item \textbf{Chương 3}: Cài đặt và thực nghiệm -- Môi trường phần cứng/phần mềm, CombatAgent, cấu hình PPO và ma trận thí nghiệm loại trừ (9 cấu hình chính + run59 + run62).
	\item \textbf{Chương 4}: Kết quả thực nghiệm -- Độ đo và tiêu chí đánh giá, sau đó phân tích kết quả nhóm theo câu hỏi nghiên cứu: kỹ thuật khám phá, curriculum learning, vùng bão hòa entropy và nút thắt quan sát (observation bottleneck).
	\item \textbf{Kết luận}: Tổng kết kết quả, hạn chế và hướng nghiên cứu tiếp theo.
\end{itemize}
